\documentclass{article}
\usepackage[utf8]{inputenc}
\usepackage[spanish]{babel}
\usepackage{listings}
\usepackage{xcolor}
\usepackage{geometry}
\geometry{a4paper, margin=1in}

\% Configuración para código C++
\definecolor{codegreen}{rgb}{0,0.6,0}        
\definecolor{codegray}{rgb}{0.5,0.5,0.5}     
\definecolor{codepurple}{rgb}{0.58,0,0.82}   
\definecolor{backcolour}{rgb}{0.95,0.95,0.92}

\lstset{
    backgroundcolor=\color{backcolour},      
    commentstyle=\color{codegreen},
    keywordstyle=\color{magenta},
    numberstyle=\tiny\color{codegray},       
    stringstyle=\color{codepurple},
    basicstyle=\ttfamily\small,
    breakatwhitespace=false,
    breaklines=true,
    captionpos=b,
    keepspaces=true,
    numbers=left,
    numbersep=5pt,
    showspaces=false,
    showstringspaces=false,
    showtabs=false,
    tabsize=2,
    language=C++,
    inputencoding=utf8,
    extendedchars=true,
    literate={á}{{\'a}}1 {é}{{\'e}}1 {í}{{\'i}}1 {ó}{{\'o}}1 {ú}{{\'u}}1 {ñ}{{\~n}}1 {¿}{{\textquestiondown}}1 {¡}{{\textexclamdown}}1
}

\% Macros para las secciones de los apuntes
\newcommand{\quees}[1]{
    \section*{🤔 ¿QUÉ ES?}
    \hrulefill \\
    \#1
}

\newcommand{\dichodeotraforma}[1]{
    \section*{💬 DICHO DE OTRA FORMA}
    \hrulefill \\
    \#1
}

\newcommand{\diagrama}[1]{
    \section*{🎨 DIAGRAMA}
    \hrulefill \\
    \#1
}

\newcommand{\comofunciona}[1]{
    \section*{⚙️ ¿CÓMO FUNCIONA?}
    \hrulefill \\
    \#1
}

\newcommand{\complejidad}[1]{
    \section*{📱 COMPLEJIDAD (Big-O)}
    \hrulefill \\
    \#1
}

\newcommand{\entrevistas}[1]{
    \section*{🎯 EN ENTREVISTAS}
    \hrulefill \\
    \#1
}

\newcommand{\resumen}[1]{
    \section*{📋 RESUMEN RÁPIDO}
    \hrulefill \\
    \#1
}

\begin{document}

\title{Merge Sort (Ordenamiento por Mezcla)}
\author{Aldo Zepeda Montoya}
\date{\today}
\maketitle

\subsection*{CATEGORÍA: 04\_Algoritmos\_Ordenamiento}

\quees{
Dividir un problema grande con un amigo 🤝. Tienes una pila gigante 
de cartas desordenadas. En lugar de ordenarlas tú solo, le das la 
mitad a tu amigo.

Cada uno ordena su mitad (y si son muchas, vuelven a dividir con 
alguien más). Al final, cuando cada quien tiene su pequeña parte 
ordenada, las fusionan (merge) comparando carta por carta para armar 
la pila final perfecta.

Es un algoritmo basado en la estrategia "Divide y Vencerás" 
(Divide and Conquer).
}

\dichodeotraforma{
Merge Sort es un algoritmo de ordenamiento recursivo, estable y de
comparación. 

Su funcionamiento se basa en dos fases:
1. Divide: Parte el array recursivamente a la mitad hasta llegar a 
   unidades de tamaño 1.
2. Conquer (Merge): Combina los subarrays ordenados para formar uno
   más grande, manteniendo el orden.

A diferencia de Quick Sort, Merge Sort GARANTIZA un tiempo de 
ejecución de O(n log n) incluso en el peor caso posible.
}

\diagrama{\begin{verbatim}
Array inicial: [38, 27, 43, 3]

Divide:
[38, 27]       [43, 3]
[38]  [27]     [43]  [3]

Merge (ordenando al combinar):
[27, 38]       [3, 43]
      \         /
      [3, 27, 38, 43]  <-- ¡Ordenado!
\end{verbatim}}

\begin{lstlisting}[language=C++]
1. Si el array tiene tamaño 1 o 0, ya está ordenado (caso base).
2. Encuentras el punto medio: `mid = size / 2`.
3. Llamas a `mergeSort` para la mitad izquierda.
4. Llamas a `mergeSort` para la mitad derecha.
5. Fucionas (Merge) las dos mitades:
   - Comparas los elementos de cada mitad uno por uno.
   - Copias el menor al array original.
   - Repites hasta terminar ambas mitades.

📝 PSEUDOCÓDIGO
──────────────────────────────────────────────────────────────

función mergeSort(arr):
    SI arr.tamaño <= 1: return arr
    
    mitad = arr.tamaño / 2
    izq = mergeSort(arr[0...mitad])
    der = mergeSort(arr[mitad...fin])
    
    return merge(izq, der)

función merge(izq, der):
    resultado = []
    MIENTRAS izq Y der NO vacíos:
        SI izq[0] <= der[0]:
            agregar izq[0] a resultado, quitar de izq
        SINO:
            agregar der[0] a resultado, quitar de der
    agregar lo que sobre de izq o der a resultado
    return resultado

💻 CÓDIGO C++
──────────────────────────────────────────────────────────────

#include <iostream>
#include <vector>
using namespace std;

void merge(vector<int>& arr, int l, int m, int r) {
    int n1 = m - l + 1;
    int n2 = r - m;
    vector<int> L(n1), R(n2);

    for (int i = 0; i < n1; i++) L[i] = arr[l + i];
    for (int j = 0; j < n2; j++) R[j] = arr[m + 1 + j];

    int i = 0, j = 0, k = l;
    while (i < n1 && j < n2) {
        if (L[i] <= R[j]) arr[k++] = L[i++];
        else arr[k++] = R[j++];
    }
    while (i < n1) arr[k++] = L[i++];
    while (j < n2) arr[k++] = R[j++];
}

void mergeSort(vector<int>& arr, int l, int r) {
    if (l >= r) return;
    int m = l + (r - l) / 2;
    mergeSort(arr, l, m);
    mergeSort(arr, m + 1, r);
    merge(arr, l, m, r);
}

int main() {
    vector<int> data = {38, 27, 43, 3, 9, 82, 10};
    mergeSort(data, 0, data.size() - 1);
    for(int x : data) cout << x << " "; // 3 9 10 27 38 43 82
    return 0;
}

⏱️ COMPLEJIDAD (Big-O)
──────────────────────────────────────────────────────────────

  Caso        │ Tiempo        │ Espacio
  ────────────┼───────────────┼───────────────
  Best        │ O(n log n)    │ O(n)
  Average     │ O(n log n)    │ O(n)
  Worst       │ O(n log n)    │ O(n)

* El espacio O(n) es porque necesita arrays auxiliares para mezclar.
\end{lstlisting}

\entrevistas{
✅ "Sort a Linked List": Merge Sort es el algoritmo preferido para 
   ordenar Linked Lists, porque no necesita acceso aleatorio.

💡 Estabilidad: Merge Sort es un Stable Sort. Esto es vital si 
   ordenas objetos por varios criterios (ej: ordenar personas por 
   nombre y luego por edad).

⚠️ Memoria: Su punto débil es el espacio O(n). Si la memoria es muy
   limitada, podrías preferir Quick Sort o Heap Sort.

🔑 Merge Step: La lógica de la función `merge` se usa en otros problemas
   como "Merge Two Sorted Lists" o "Merge Intervals".
}

\resumen{
• Divide y vencerás.
• Siempre tarda O(n log n), sin importar qué tan desordenado esté.
• Es un ordenamiento estable.
• Usa memoria extra O(n).
• Ideal para datasets grandes o Linked Lists.
• Base para problemas de fusionar listas o intervalos.
════════════════════════════════════════════════════════════════
}

\end{document}